\section{Introduction to Optimization}

Every optimization problem asks the same question: among all available choices, which one is best?
The answer depends on what "best" means — a cost to minimize, a reward to maximize, a trade-off to balance — and on what choices are actually available.
This chapter sets up the vocabulary and the mathematical objects that make that question precise.


\subsection{Optimization Problems}

\begin{defn}[Univariate function]
  A function $f \colon \R \to \R$ maps an input $x$ to an output $y = f(x)$.
  Four sets capture its geometry:
  \begin{itemize}
    \item $\mathrm{Gr}(f) = \{(x, f(x)) \mid x \in \R\} \subset \R^2$
      \hfill\textit{[Graph]}
    \item $\mathrm{Im}(f) = \{y \mid \exists\, x : y = f(x)\} \subset \R$
      \hfill\textit{[Image]}
    \item $L(f, v) = \{x \mid f(x) = v\} \subset \R$
      \hfill\textit{[Level set at value $v$]}
    \item $L(f, 0) = \{x \mid f(x) = 0\} \subset \R$
      \hfill\textit{[Roots]}
  \end{itemize}
\end{defn}

\begin{defn}[Minimization problem]\label{def:minimization}
  Given a function $f$ and a set $X$, we seek $x^*$ that minimizes $f$ over $X$.
  Two cases arise:
  \begin{itemize}
    \item \textbf{Unconstrained:}
      $(P)\; f^* = \min\{f(x) \mid x \in \R\}$,\;
      $x^* \in \argmin\{f(x) \mid x \in \R\}$
    \item \textbf{Constrained:}
      $(P)\; f^* = \min\{f(x) \mid x \in X\}$,\;
      $x^* \in \argmin\{f(x) \mid x \in X\}$
  \end{itemize}
  The minimizer $x^*$ need not be unique, but every minimizer yields the same optimal value $f^*$.
  For constrained problems we call $X$ the \textit{feasible region},
  any $x \in X$ a \textit{feasible solution},
  any $x \notin X$ an \textit{infeasible solution},
  and $X^* = L(f, f^*) \cap X$ the \textit{optimal set}.
\end{defn}

\begin{defn}[Types of constraints]
  Constraints that describe $X$ fall into three families:
  \begin{itemize}
    \item $g(x) = v$ \hfill\textit{[Equality — a level set]}
    \item $g(x) \leq v$ \hfill\textit{[Inequality — a sublevel set]}
    \item $g_1(x) \leq 0 \;\land\; g_2(x) \leq 0 \;\land\; \cdots \;\land\; g_k(x) \leq 0$
      \hfill\textit{[Multiple — an intersection]}
  \end{itemize}
\end{defn}

\begin{defn}[Function unbounded below]
  A function $f$ is \textit{unbounded below} if for every $M > 0$ there exists some $x$ with $f(x) < -M$.
  No finite minimum exists, and the infimum is $f^* = -\infty$.
\end{defn}


\subsection{Solution Quality}

Exact minimizers are rarely available in practice.
Algorithms produce sequences of iterates that get closer and closer to $f^*$, and we need a way to say "close enough."

\begin{defn}[Approximate solution]
  Fix a tolerance $\varepsilon > 0$.
  A point $x_\varepsilon$ is an \textit{$\varepsilon$-optimal solution} of $(P)$ if
  $$
    f^* \leq f(x_\varepsilon) \leq f^* + \varepsilon.
  $$
  How hard it is to find $x_\varepsilon$ depends on $\varepsilon$, on $f$, and on $X$.
\end{defn}

Two standard measures quantify the gap between a candidate $x$ and the true optimum:
\begin{itemize}
  \item $A(x) = f(x) - f^*$
    \hfill\textit{[Absolute gap]}
  \item $R(x) = A(x) \,/\, \abs{f^*}$
    \hfill\textit{[Relative gap]}
\end{itemize}
Both require knowledge of $f^*$, which is typically unknown.
Much of optimization theory is devoted to bounding these gaps \emph{without} knowing $f^*$ — through duality, descent lemmas, or convergence-rate analysis.


\subsection{Multivariate and Multi-objective Problems}

\begin{defn}[Multivariate optimization]
  When $f \colon \R^n \to \R$ with $n > 1$, the feasible region can take many forms.
  A simple but common one is the \textit{box}:
  $$
    X = \{x \in \R^n : x_- \leq x \leq x_+\},
    \qquad x_-, x_+ \in \R^n,\; x_- \leq x_+,
  $$
  where the inequalities hold component-wise.
\end{defn}

What happens when we have more than one objective?

\begin{defn}[Multi-objective optimization]\label{def:multi-objective}
  Given $f \colon \R^n \to \R^k$ with $k > 1$, the problem becomes
  $$
    (P) \quad \min\bigl\{[f_1(x),\, f_2(x),\, \dots,\, f_k(x)] : x \in X\bigr\}.
  $$
  Since $\R^k$ for $k > 1$ lacks a total order, there is no single "best" solution.
  We need a weaker notion of optimality.
\end{defn}

\begin{defn}[Pareto optimality]\label{def:pareto}
  A feasible point $x^*$ is \textit{Pareto-optimal} if no other $x \in X$ satisfies
  $f_i(x) \leq f_i(x^*)$ for all $i$ with strict inequality for at least one $i$.
  The set of all Pareto-optimal points is the \textit{Pareto frontier}.
\end{defn}

The Pareto frontier is typically a surface (or curve, when $k = 2$), and a decision-maker must choose a point on it.
Two common strategies reduce the multi-objective problem to a single-objective one:
\begin{itemize}
  \item $\min\{f_1(x) + \alpha\, f_2(x) : x \in X\}$
    \hfill\textit{[Scalarization]}
  \item $\min\{f_1(x) : f_2(x) \leq \beta,\; x \in X\}$
    \hfill\textit{[Constraint-based budgeting]}
\end{itemize}
Different values of $\alpha$ or $\beta$ trace out different points on the frontier.


% ────────────────────────────────────────────────────────────────
\section{Spectral Theory and Quadratic Forms}

Quadratic functions appear everywhere in optimization — as objectives, as local approximations of smooth functions, and as building blocks of convergence proofs.
Understanding them requires a detour through eigenvalues and matrix structure.
The payoff: once we know the spectrum of a matrix, we can predict almost everything about the associated quadratic.


\subsection{Tomographies}

How do we study a function of $n$ variables?
One powerful strategy is to "slice" it along a direction and analyze the resulting one-dimensional function.

\begin{defn}[Tomography]\label{def:tomography}
  Given $f \colon \R^n \to \R$, an origin $x \in \R^n$, and a direction $d \in \R^n$, the \textit{tomography} of $f$ along $d$ from $x$ is
  $$
    \varphi_{x,d}(\alpha) = f(x + \alpha d) \colon \R \to \R.
  $$
\end{defn}

A tomography reduces a multivariate problem to a univariate one.
Line search methods exploit this directly: pick a descent direction $d$, then minimize $\varphi_{x,d}$ over $\alpha \geq 0$.
We will see this idea at work in every gradient-based algorithm.


\subsection{Eigenvalues and Spectral Decomposition}

\begin{defn}[Eigenvalues and eigenvectors]\label{def:eigen}
  Given a matrix $Q \in \R^{n \times n}$, a scalar $\lambda \in \R$ is an \textit{eigenvalue} and a nonzero vector $v \in \R^n$ is the corresponding \textit{eigenvector} if
  $$
    Qv = \lambda v.
  $$
  Equivalently, $v$ is a direction along which $Q$ acts as a simple rescaling.
\end{defn}

Eigenvalues can be complex in general, and eigenvectors need not be orthogonal.
Symmetric matrices are the exception — and the one that matters most for optimization.

\begin{thm}[Eigenvalues of symmetric matrices]\label{thm:sym-eigen}
  If $Q \in \R^{n \times n}$ is symmetric, then $Q$ has $n$ real eigenvalues (counted with multiplicity) and $n$ eigenvectors that can be chosen orthonormal.
\end{thm}

We denote the $n$ orthonormal eigenvectors as $H_1, H_2, \dots, H_n \in \R^n$ with corresponding eigenvalues $\lambda_1, \lambda_2, \dots, \lambda_n \in \R$, so that
$$
  Q H_i = \lambda_i H_i \quad \forall\, i \in \{1, 2, \dots, n\}.
$$
We collect these eigenvectors into a single matrix.

\begin{defn}[Eigenvector matrix]\label{def:eigvec-matrix}
  For a symmetric $Q$ with orthonormal eigenvectors $H_1, H_2, \dots, H_n \in \R^n$ and corresponding eigenvalues $\lambda_1 \geq \lambda_2 \geq \cdots \geq \lambda_n$, define
  $$
    H = [H_1 \mid H_2 \mid \cdots \mid H_n] \in \R^{n \times n}.
  $$
  Orthonormality of the columns means $H^T H = I$, so $H$ is an orthogonal matrix.
\end{defn}

\begin{thm}[Spectral decomposition]\label{thm:spectral}
  Every symmetric $Q \in \R^{n \times n}$ can be written as
  $$
    Q = H \Lambda H^T = \sum_{i=1}^{n} \lambda_i\, H_i H_i^T,
  $$
  where $\Lambda = \mathrm{diag}(\lambda_1, \dots, \lambda_n)$.
  Each term $\lambda_i\, H_i H_i^T$ is a rank-one matrix: the decomposition expresses $Q$ as a weighted sum of $n$ such terms, one per eigenvalue.
\end{thm}

Why does this matter?
The spectral decomposition turns matrix operations into scalar ones.
If we change coordinates via $y = H^T x$, the quadratic $x^T Q x$ becomes $\sum_i \lambda_i y_i^2$ — a sum of independent squared terms, each scaled by an eigenvalue.
Large eigenvalues stretch, small ones compress.
The \emph{ratio} between the largest and smallest controls how "distorted" the quadratic is, and (as we will see in the next chapter) how fast gradient methods converge.

\begin{thm}[Variational characterization of eigenvalues]\label{thm:var-eigen}
  For a symmetric $Q$, the extreme eigenvalues satisfy
  $$
    \lambda_1 = \max\left\{\frac{d^T Q d}{d^T d} : d \neq 0\right\},
    \qquad
    \lambda_n = \min\left\{\frac{d^T Q d}{d^T d} : d \neq 0\right\}.
  $$
  The ratio $d^T Q d \,/\, d^T d$ is the \textit{Rayleigh quotient}.
  The maximum is attained at $d = H_1$, the minimum at $d = H_n$.
\end{thm}


\subsection{Matrix Definiteness}

\begin{defn}[Determinant]\label{def:determinant}
  The determinant of $Q$ equals the product of its eigenvalues:
  $$
    \det(Q) = \prod_{i=1}^{n} \lambda_i.
  $$
  If $\det(Q) \neq 0$, then $Q$ is nonsingular, every eigenvalue is nonzero, and $Q^{-1}$ exists.
  The eigenvectors of $Q^{-1}$ are the same as those of $Q$, but with reciprocal eigenvalues: if $Qv = \lambda v$ then $Q^{-1} v = \frac{1}{\lambda}\, v$.
  It follows that $\det(Q^{-1}) = \det(Q)^{-1}$.
\end{defn}

The sign pattern of the eigenvalues determines the shape of the associated quadratic form — and therefore the nature of its critical points.

\begin{defn}[Matrix definiteness]\label{def:definiteness}
  A symmetric matrix $Q$ is:
  \begin{itemize}
    \item \textit{Positive definite} ($Q \succ 0$): all $\lambda_i > 0$.
      The quadratic $\frac{1}{2} x^T Q x$ has a unique global minimum at the origin.
    \item \textit{Negative definite} ($Q \prec 0$): all $\lambda_i < 0$.
      Unique global maximum at the origin.
    \item \textit{Positive semidefinite} ($Q \succeq 0$): all $\lambda_i \geq 0$.
      The quadratic is convex, but the minimum (if it exists) need not be unique — flat directions correspond to zero eigenvalues.
    \item \textit{Negative semidefinite} ($Q \preceq 0$): all $\lambda_i \leq 0$.
      Concave, with a possibly non-unique maximum.
    \item \textit{Indefinite}: eigenvalues of both signs.
      The origin is a saddle point.
  \end{itemize}
\end{defn}


\subsection{Quadratic Forms: Analysis and Examples}

The general quadratic function in $\R^n$ is
$$
  f(x) = \frac{1}{2} x^T Q x + \inner{q, x} + c,
  \qquad Q \in \R^{n \times n},\; q \in \R^n,\; c \in \R.
$$
We analyze two cases: the homogeneous one ($q = 0$, $c = 0$) and the general one.

\begin{exm}[Homogeneous quadratic form]
  Set $q = 0$ and $c = 0$.
  Then
  $$
    f(x) = \frac{1}{2} x^T Q x
    = \frac{1}{2}\biggl[\sum_{i=1}^{n} Q_{ii}\, x_i^2
      + \sum_{\substack{i,j=1 \\ j \neq i}}^{n} Q_{ij}\, x_i x_j\biggr].
  $$
  An important observation: even if $Q$ is not symmetric, we can always replace it with the symmetric matrix $(Q + Q^T)/2$ without changing $f$.
  Indeed,
  $$
    x^T Q x
    = \frac{(x^T Q x) + (x^T Q x)^T}{2}
    = x^T \frac{Q + Q^T}{2}\, x.
  $$
  So there is no loss of generality in assuming $Q$ symmetric from the start.

  The tomography along a direction $d$ from the origin gives
  $$
    \varphi(\alpha) = f(\alpha d) = \frac{1}{2}\, \alpha^2 (d^T Q d).
  $$
  This is a univariate parabola whose sign depends entirely on $d^T Q d$.
  When $Q \succ 0$, the parabola opens upward for every $d \neq 0$ — the function has a unique minimum at the origin.
  When $Q$ is indefinite, some directions yield upward parabolas and others downward: a saddle.
\end{exm}

\begin{exm}[Non-homogeneous quadratic function]
  Consider the full quadratic
  $$
    f(x) = \frac{1}{2}\, x^T Q x + \inner{q, x} + c.
  $$
  The linear term $\inner{q, x}$ breaks the symmetry around the origin.
  Can we get rid of it?

  \textbf{Case 1 — $Q$ nonsingular.}
  Setting $\nabla f(x) = Qx + q = 0$ yields the critical point $\bar{x} = -Q^{-1} q$.
  Think of $\bar{x}$ as the point where the quadratic "wants" to settle — the natural center of the function.
  We shift coordinates to that center by defining $z = x - \bar{x}$.
  Substituting $x = z + \bar{x}$ and using $Q \bar{x} = -q$, every linear term cancels:
  $$
    f(x) = \frac{1}{2}\, z^T Q z + f(\bar{x}).
  $$
  The problem reduces to a homogeneous quadratic in $z$ plus a constant offset.
  If $Q \succ 0$, the unique minimum sits at $z = 0$, i.e.\ $x = \bar{x}$.
  If $Q$ is indefinite, $\bar{x}$ is a saddle point.

  \textbf{Case 2 — $Q$ singular.}
  When $Q$ is singular, $Q^{-1}$ doesn't exist and the trick above fails.
  We decompose the linear term as $q = q^0 + q^+$, where
  $$
    q^0 \in \ker(Q), \qquad q^+ \in \mathrm{im}(Q).
  $$
  The component $q^+$ lives in the range of $Q$, so we can still "absorb" it via a partial change of variables.
  What matters is $q^0$, the part that projects onto the null space:
  \begin{itemize}
    \item If $q^0 = 0$: a minimum exists, but it is not unique.
      The entire affine subspace $\bar{x} + \ker(Q)$ consists of minimizers — the function is flat along null-space directions.
    \item If $q^0 \neq 0$: no minimum exists.
      Moving along the null-space direction $q^0$, the quadratic term $\frac{1}{2} x^T Q x$ stays constant (because $Q$ kills that direction), but the linear term $\inner{q^0, x}$ keeps decreasing.
      The function is unbounded below.
  \end{itemize}
\end{exm}